\documentclass[11pt]{article}
\usepackage[utf8]{inputenc}
\usepackage{amsmath,amssymb}
\usepackage{authblk}
\usepackage[margin=1in]{geometry}
\usepackage[hidelinks]{hyperref}
\usepackage{xurl}
\usepackage{setspace}

\title{Energy Without Mass: RealQM, Inertia, and the Penning Trap}
\author[1]{Claes Johnson}
\affil[1]{KTH Royal Institute of Technology, SE-100 44 Stockholm, Sweden \\
\texttt{claesjohnson@gmail.com}}
\date{\today \\[0.5em]
\small\textit{Written in respectful cooperation between Claes Johnson and Claude (Anthropic), with Claude
contributing a balanced, neutral view.}}

\begin{document}
\maketitle

\begin{abstract}
RealQM computes the energy of matter --- every atom, and nuclear binding --- from one dimensionless functional:
a localisation cost plus the Coulomb energy of charge densities in ordinary three-dimensional space. \emph{Mass
never enters this computation.} The coefficient of the localisation term, written $\hbar^2/2m$ and read in
textbooks as a kinetic energy $p^2/2m$, is here only a geometric scale --- a stationary state has no motion, so it
fixes a length and an energy, not a mass --- and atomic energies come out as pure numbers in Hartree, with the
electron mass merely the unit. If mass and energy are equivalent, one of the two words is redundant; we keep
energy. The apparent obstacle is the Penning trap, which seems to \emph{weigh}. We model it in energy alone: an ion
of energy $E$ has inertia $E/c^2$ --- the momentum its own field carries when it moves --- and circles at
$\omega_c=qBc^2/E$, reproducing the electron and proton cyclotron frequencies and rendering the He\,-4 ``mass
defect'' as the binding energy read as a frequency shift. The one relation this needs, $\mathbf p=(E/c^2)\mathbf v$,
is the measured law of inertia; RealQM's part is to show why it is natural --- the localisation term is a genuine
cohesive stress, so every bound state is a stress-balanced equilibrium, and indeed the time-independent
Schr\"odinger equation \emph{is} this local force balance. The account lives in a single inertial frame with
Maxwell's equations and the measured inertia law; a first-principles derivation of the coefficient (disposing of
the historical $4/3$) lies beyond this elementary framework but is not needed here. Mass is thereby an output, not an input; the elementary rest energies remain, as in every theory, the
inputs no theory derives.
\end{abstract}

\setstretch{1.06}

\paragraph{What is new, and what is standard.} The physics this note relies on is standard: mass--energy
equivalence and $E/c^2$ inertia, the electromagnetic field momentum and its historical factor $4/3$ (Thomson,
Abraham), the Penning trap as a measurement of inertia, and the hydrodynamic reading of the stationary
Schr\"odinger equation (Madelung, Bohm). No new empirical fact and no new mathematics is claimed. What is the
author's is the framework --- RealQM, quantum mechanics as real charge densities in three-space --- and, within
it, a synthesis: that mass is not a primitive but the energy content read as inertia; that the localisation term
already \emph{is} the cohesive stress classical electrodynamics had to postulate; and that the whole can be told in
one frame. The note is foundational and conceptual, not a report of new physics.

\section{The kinetic coefficient is a geometric scale, not a mass}
The energy RealQM~\cite{RealQM} minimises for $N$ electrons is
\begin{equation}
E(\psi)=\tfrac12\int|\nabla\psi|^2
        -\sum_a\int\frac{Z_a}{|\mathbf x-\mathbf X_a|}\,\rho
        +\tfrac12\iint\frac{\rho(\mathbf x)\rho(\mathbf x')}{|\mathbf x-\mathbf x'|},
\qquad \rho=\psi^2,
\label{eq:functional}
\end{equation}
a localisation cost plus Coulomb attraction to the kernels and Coulomb repulsion between the clouds, with no
self-interaction and non-overlapping supports. In physical units the first term is $(\hbar^2/2m)\int|\nabla\psi|^2$,
and standard quantum mechanics reads it as the kinetic energy $p^2/2m$ of a particle in motion. But in a
stationary state nothing moves: $\langle\mathbf p\rangle=0$ and the density is static. The ``velocity'' has no
referent, and $p^2/2m$ is a formality. What the term physically is, is a \emph{cost of localising static charge}:
the finer the density's variation, the larger $\int|\nabla\psi|^2$. Its coefficient does one job --- it fixes,
against the Coulomb terms, the characteristic \emph{size} and \emph{energy} of the minimiser. It is a geometric
scale, and the symbol $m$ in it is inherited machinery, not a mass that the theory reasons about.

\section{Atomic energies without mass}
In atomic units ($\hbar=m_e=e=1$) the functional \eqref{eq:functional} is dimensionless and its minimum is a pure
number: $-0.5$ for H, $-2.90$ for He, $-7.48$ for Li, and so on across the periodic table. The whole computation ---
a Laplacian, a $1/r$, non-overlap, downhill relaxation --- is the geometry of charge in three-space; no mass
appears as a quantity anywhere in it. Mass enters in exactly one place, the overall unit
$1~\mathrm{Hartree}=m_e e^4/\hbar^2 = 27.2~\mathrm{eV}$, a single universal constant, identical for every atom,
light or heavy --- the ruler, not an ingredient. Ratios, trends, and the entire periodic structure are mass-free.

This is sharpest for heavy atoms, where the textbook invokes a velocity-dependent mass increase for fast inner
electrons. RealQM uses the \emph{identical} functional for uranium as for hydrogen: the electrons are static, so
there is no velocity to invoke and no mass to increase; the residual few percent is the spherically-symmetric
simplification. One mass-free functional covers the whole table.

\section{The nucleus: geometry gives the ratios, a scale gives the rest}
Model a nucleus the same way --- charge clouds minimising \eqref{eq:functional} on nested non-overlapping shells.
The dimensionless \emph{structure} is again mass-free, and RealQM reproduces it: the He\,-4/deuteron binding ratio
comes out $\approx 13$ against an experimental $\approx 12.7$, from Coulomb geometry alone. The \emph{absolute}
scale --- why nuclear binding is measured in MeV and lengths in fermis rather than eV and \aa{}ngstr\"oms --- is
set by the coefficient of the localisation term: a smaller coefficient contracts the clouds and deepens the
binding, exactly as rescaling the muonic hydrogen curve does. It is tempting to read this coefficient as a mass
and then to require a large ``glue mass'' to reach the nuclear scale. That requirement is an artifact of the mass
reading. The coefficient is a geometric scale; the nuclear scale is one number, fixed by matching a single datum
(a nuclear length, or one binding energy), after which the mass-free ratios follow. No glue mass need be sought.

\section{Mass is not primitive; it is energy showing up as inertia}
If mass and energy are genuinely equivalent, then ``mass'' is a redundant word and the physics can be told with
energy alone. This is not a mere convention. A Penning trap measures a real thing --- \emph{inertia}, a bound
system's resistance to acceleration --- and the inertia deficit of a bound nucleus (He\,-4 is $0.76\%$ lighter
than its constituents) is a measured mechanical fact that coincides, through $c^2$, with the binding energy
measured independently by reaction $Q$-values (quantified in \S6). The lesson is not two things linked by a
formula, but \emph{one} thing: energy, manifesting dynamically as inertia. ``Mass'' names that inertia --- an
output, never an input, droppable without loss.

\section{A RealQM model of the Penning trap}
The trap is precisely a machine for measuring inertia: a charge $q$ in a uniform field $B$ circles at the
cyclotron frequency $\omega_c=qB/m$. To model it without mass, the inertia must be energy. Two things must be kept
apart here. Maxwell supplies the \emph{mechanism}: a moving charge drags its own field, and that field carries
momentum, $\mathbf p\propto(\text{energy}/c^2)\,\mathbf v$ (Thomson's electromagnetic mass), so inertia is field
momentum, not a separate substance. But Maxwell does not supply the exact \emph{coefficient}: the bare field
integral gives the classical factor $4/3$ and only the field energy, not the total. The exact relation
\begin{equation}
\mathbf p=\frac{E}{c^2}\,\mathbf v
\label{eq:pev}
\end{equation}
--- unit coefficient, total energy --- is the \emph{measured} law (\S4, and \S6 for its status): the trap gives
the inertia $qB/\omega_c$, and it is found to equal $E/c^2$ to a part in $10^{7}$. So the field momentum is the
mechanism and the measurement fixes the coefficient; with \eqref{eq:pev} the cyclotron frequency is then
\begin{equation}
\omega_c=\frac{qB}{E/c^2}=\frac{qBc^2}{E},
\label{eq:wc}
\end{equation}
written in charge, energy, $B$ and $c$, with no mass symbol. Evaluated at $B=1$~T:
\begin{center}
\begin{tabular}{lll}
\hline
system & energy $E$ & $\omega_c/2\pi$ from \eqref{eq:wc} \\
\hline
electron & $0.511$ MeV & $27.99$ GHz/T \quad(measured $27.99$)\\
proton   & $938.27$ MeV & $15.24$ MHz/T \quad(measured $\sim15.2$)\\
He\,-4 (free $\to$ bound) & $-28.30$ MeV binding & $\omega_c$ shifts $+0.76\%$ \\
\hline
\end{tabular}
\end{center}
The electron and proton cyclotron frequencies come out of their energies alone. And the nuclear ``mass defect''
is nothing but the binding energy read as a frequency shift: bound He\,-4 has $28.3$~MeV less energy, hence a
$0.76\%$ higher $\omega_c$ (smaller $E$ in the denominator of \eqref{eq:wc}). The trap does not weigh a substance;
it reports an energy. Mass has been eliminated.

\section{The bound state is a stress-balanced equilibrium}
The trap model uses one relation, $\mathbf p=(E/c^2)\mathbf v$ --- that a moving cloud's inertia is exactly its
energy. This is the \emph{measured} law of \S4 and needs no derivation; what RealQM adds is the reason it is
\emph{natural}. Computed from Maxwell, the field-\emph{alone} momentum of a moving charge is the classical value
$\tfrac43(U_{\text{field}}/c^2)\mathbf v$~\cite{Thomson,Abraham} --- not the total, and no paradox: the remainder is
carried by whatever stress holds the charge together, the ``Poincar\'e stress'' classical electrodynamics had to
\emph{postulate}. RealQM need not postulate it. The localisation term $\tfrac12\int|\nabla\psi|^2$ \emph{is} that
cohesive stress, and a bound state is by construction a stress-balanced equilibrium: as the minimiser of $E=T+V$ it
is stationary under the scaling $\psi(\mathbf x)\mapsto\lambda^{3/2}\psi(\lambda\mathbf x)$, giving the virial
identity
\begin{equation}
2T+V=0,
\label{eq:virial}
\end{equation}
so the internal stresses balance,
\begin{equation}
\int T^{ij}\,d^3x=0 \qquad\text{(mechanical equilibrium).}
\label{eq:laue}
\end{equation}
That every RealQM bound cloud is such an equilibrium is a single-frame fact, and it is the whole of RealQM's
contribution here: the cohesive stress classical electrodynamics had to \emph{invent} is, in RealQM, the
localisation energy it already has. The total inertia is then the measured $E_{\text{total}}/c^2$, the field-alone
$4/3$ being merely a part of the whole. A first-principles \emph{derivation} of the exact coefficient, disposing of
the $4/3$, lies beyond this elementary framework and is not pursued here; the coefficient is taken from measurement
(\S4), and the single-frame account needs no more.

That inertia \emph{is} energy is, moreover, no definition smuggled in but a tested identity between two independent
measurements. A Penning trap returns a cyclotron frequency --- an $F/a$, pure inertia --- from which one reads the
mass difference between a nucleus and its separated constituents; the \emph{same} binding energy is measured, with
no trap and no mass, by calorimetry of the reaction --- the $\gamma$-ray energy $Q$ released when the constituents
combine. The two agree, $\Delta m\,c^2=Q$, to about one part in $10^{7}$~\cite{Rainville}. A frequency on one side, a photon energy
on the other, joined only by $c^2$: it is this agreement, not a choice of words, that makes ``inertia is energy''
a physical statement. RealQM's part is to compute the difference; that the inertia equals it is the measured law of \S4.

\section{The stress tensor made explicit: the localisation term as the cohesive stress}
The settlement of the previous section can be shown directly, at the level of the stress tensor, with no appeal to
a postulated cohesion. A RealQM cloud is a real field $\psi$, $\rho=\psi^2$, obeying the stationarity condition
(atomic units)
\begin{equation}
-\tfrac12\nabla^2\psi + U\psi = \varepsilon\psi
\qquad\Longleftrightarrow\qquad
\nabla^2\psi = 2(U-\varepsilon)\psi,
\label{eq:stat}
\end{equation}
with $U$ the potential the cloud sees --- external plus the Coulomb potential of the \emph{other} charges. The
momentum flux carried by the localisation term $\tfrac12|\nabla\psi|^2$ is the matter stress tensor
\begin{equation}
T^{M}_{ij} = \partial_i\psi\,\partial_j\psi - \delta_{ij}\Big[\tfrac12|\nabla\psi|^2 + (U-\varepsilon)\rho\Big].
\label{eq:TM}
\end{equation}
Taking its divergence and using \eqref{eq:stat} together with
$\partial_i(\partial_i\psi\,\partial_j\psi)=(\nabla^2\psi)\,\partial_j\psi+\tfrac12\partial_j|\nabla\psi|^2$, every
term collapses to one:
\begin{equation}
\boxed{\;\partial_i T^{M}_{ij} = -\rho\,\partial_j U\;}
\label{eq:balance}
\end{equation}
verified symbolically and, on a three-dimensional grid for the hydrogen ground state, numerically to the
finite-difference error. Equation~\eqref{eq:balance} is local Newton's second law for the static cloud: the
divergence of the \emph{localisation} stress equals minus the force density $\rho\,\nabla U$. The quantum pressure
supplies precisely the stress that balances the Coulomb force --- the cohesive ``Poincar\'e stress'' is here a
\emph{consequence} of the stationary equation, not an addition to it.

\paragraph{The stress balance is the time-independent Schr\"odinger equation.}
Equation~\eqref{eq:balance} is not merely \emph{implied} by \eqref{eq:stat}; for the real, nodeless ground state
it is equivalent to it, and the equivalence exposes what Schr\"odinger's stationary equation \emph{is}. Writing
$\psi=\sqrt\rho$ and dividing \eqref{eq:stat} by $\psi$ gives
\begin{equation}
U + Q = \varepsilon,\qquad Q \equiv -\tfrac12\,\frac{\nabla^2\sqrt\rho}{\sqrt\rho},
\label{eq:madelung}
\end{equation}
hence $\nabla(U+Q)=0$: at every point the potential force $-\rho\,\nabla U$ is cancelled by the quantum-pressure
force $-\rho\,\nabla Q=\nabla\!\cdot\sigma^{Q}$, whose stress $\sigma^{Q}$ is the gradient part of \eqref{eq:TM}.
The time-independent Schr\"odinger equation is thus a statement of \emph{local mechanical equilibrium} --- Newton's
second law for a charge fluid carrying a quantum pressure --- with the eigenvalue $\varepsilon$ the constant of
integration, the uniform level (a chemical potential) at which the balance settles; no probability enters, it is
$\textstyle\sum\text{forces}=0$ for real charge. This is the Madelung--Bohm hydrodynamic form of the
equation~\cite{Madelung,Bohm}; what RealQM adds is to read the fluid as the electron itself, so that the balance is
a real equilibrium and the localisation term a real stress. The complex phase set aside here is precisely the
\emph{flow}: the time-dependent equation is this same balance plus continuity and an Euler equation for
$\mathbf v=\nabla S$ --- being's equilibrium acquiring becoming's dynamics.

Splitting $U=V_{\text{ext}}+\Phi$, the Coulomb piece $-\rho\,\partial_j\Phi=\rho E_j$ is cancelled by the Maxwell
stress ($\partial_i T^{\mathrm{EM}}_{ij}=\rho E_j$, the mutual fields pairing by Newton's third law), leaving
\begin{equation}
\partial_i\big(T^{M}+T^{\mathrm{EM}}\big)_{ij} = -\rho\,\partial_j V_{\text{ext}}.
\label{eq:total}
\end{equation}
For a self-bound system --- an atom, molecule or nucleus with \emph{all} its constituents included, so
$V_{\text{ext}}=0$ --- the total stress is divergence-free, and the tensor-virial identity
$\int T^{ij}\,d^3x=\int\partial_k(x^i T^{kj})\,d^3x-\int x^i\,\partial_k T^{kj}\,d^3x=0$ (the surface term vanishing
for a localised system) gives, on boosting, $\mathbf p=(E_{\text{total}}/c^2)\mathbf v$ with no residual $4/3$.
For a spherical atom, isotropy reduces \eqref{eq:total} to its trace, the virial $2T+V=0$ of \eqref{eq:virial}.

Three corollaries follow, and they are consequences of \eqref{eq:balance}, not caveats appended to it. First, the
inter-particle (binding) stresses close exactly, so the mass \emph{defect} a Penning trap measures obeys
$\mathbf p=(E/c^2)\mathbf v$ rigorously: RealQM is exact where it speaks. Second, RealQM carries no
self-interaction, so the self-field terms are simply absent from \eqref{eq:balance}; the rest inertia of a
\emph{lone} charge is not produced, and remains the elementary input discussed below. Third, a single cloud held
only by $V_{\text{ext}}$ leaves the residual $-\rho\,\partial_j V_{\text{ext}}$ of \eqref{eq:total}, which closes
only when the confiner is itself part of the system --- the trap that makes a single trapped charge a bound
\emph{geonium}. The honest boundaries of the construction are thus theorems of the same identity that settles it.

\section{One frame is enough}
The entire argument is carried out in a single frame. In it a moving charge drags a magnetic field
$\mathbf B\propto\mathbf v/c$ and its field carries the Poynting momentum $\mathbf g=\mathbf S/c^2$ --- Maxwell,
and J.~J.~Thomson's electromagnetic mass of 1881~\cite{Thomson} --- while the relation
$\mathbf p=(E/c^2)\mathbf v$ is the measured law of \S4. The account uses only this: one frame, Newton's
$\mathbf F=\dot{\mathbf p}$, Maxwell's field, and the measured inertia. Nothing further is invoked.

The historical $4/3$ is, in the one frame, no paradox. It is merely the field-\emph{alone} momentum of a moving
charge, $(4/3)(U_{\text{field}}/c^2)\mathbf v$, a legitimate quantity; the \emph{total} inertia is the measured
$E/c^2$, field and localisation stress together. As RealQM keeps Coulomb and drops Born's probabilities, this
account keeps Maxwell and the measured inertia law, and needs no further apparatus.

Two honest limits remain, and they mark the edge of this elementary framework. First, RealQM's matter description
and the Maxwell field are joined only to leading order in $v/c$ --- ample for the inertia of bound matter and for
the Penning trap, but a full account of fast-moving matter and its fine anomalies lies beyond it. Second, a
first-principles \emph{derivation} of the coefficient in $\mathbf p=(E/c^2)\mathbf v$ --- here taken from
measurement --- lies beyond the framework as well; and the absolute rest energies remain inputs, the mass-spectrum
problem, not a debt of this account. These edges are noted, not pursued.

\section{Absolute masses are inputs, not a debt}
Nothing above requires the absolute rest energies of the free constituents. Every mass measurement is a
\emph{ratio} to a reference (the atomic mass unit is $^{12}\mathrm C/12$ by definition), so what is physical is
mass ratios and differences --- and those are exactly what a Penning trap reports and what RealQM computes:
binding energies mass-free (\S2--\S3) and bound-system inertia exactly ($\S6$). The one genuinely physical
remainder is the set of dimensionless elementary ratios such as $m_e/m_p\approx 1/1836$; but deriving the mass
\emph{spectrum} from first principles is the deepest open problem of all physics --- the Standard Model inputs its
Yukawa couplings, QCD its scale --- so RealQM taking the elementary masses as inputs is universal practice, not a
debt peculiar to it. There is no ``glue mass'' to find and no absolute self-energy that any theory computes.

A natural objection sharpens this line. A RealQM charge is a cloud that needs a confining agent to be localised at
all --- a free electron, held only by its localisation term, simply spreads --- so how can a \emph{single}
electron or proton be trapped? Because in the trap it is not free: the magnetic and electrostatic fields confine
it, and Dehmelt, who first did it, named the result \emph{geonium}~\cite{Dehmelt}, an artificial atom bound to the apparatus. The
trap is the ``surrounding.'' But what it reads for a lone elementary charge is that charge's rest energy, and that
RealQM does not derive: a cloud spread over a macroscopic orbit has negligible self-field energy, so the electron's
$0.511$~MeV is the elementary baseline no theory produces, not a RealQM self-energy. The division is clean ---
RealQM's mechanism, inertia as \emph{stored, bound} energy, is what a trap confirms when it weighs a composite
\emph{defect}; a single trapped particle merely reads the input baseline. The trap thus tests RealQM exactly where
RealQM speaks (binding differences) and stays silent exactly where RealQM is silent (the elementary rest energies).

\section{Conclusion}
RealQM eliminates mass from the physics it computes. Atomic and nuclear energies are geometry --- charge and a
scale coefficient, the electron mass only a unit; the differences that every experiment reports are mass-free; and
the inertia a Penning trap weighs is, for any bound system, its energy over $c^2$ --- the measured law, with the
localisation term playing the cohesive stress that classical electrodynamics had to postulate. In this program
$E=mc^2$ is not borrowed but \emph{earned} for bound matter, and the one instrument that seems to force mass upon
us is modelled in energy alone --- in a single frame, with Maxwell's field and the measured inertia, and nothing
more. The elementary masses remain inputs, as they are in every theory --- which is to say the
concept of mass, everywhere it is actually used or measured, has been shown dispensable: charge and energy
suffice.

A last distinction places the whole result. Energy alone suffices for \emph{being}: the static structure of matter
--- atoms, molecules, nuclei, and every energy difference an experiment reports --- follows from minimising a
Coulomb functional in which no mass, no inertia and no $c$ appear. Inertia enters only for \emph{becoming}: a charge
that travels and is accelerated must be told how it resists force, and that resistance is $E/c^2$, the momentum its
field carries, drawn from Maxwell --- not a second quantity but the same energy in dynamical clothes. A magnetic
field makes this transparent, and shows the right word is \emph{charge}, not electron: the force
$q\mathbf v\times\mathbf B$ couples to the charge, the curvature $r=(E/c^2)\,v/qB$ to the energy, and neither is a
mass --- a spectrometer sorts \emph{any} charge alike by $q$ and $E/c^2$, the growth of the
curvature-inertia with speed being just the kinetic energy joining the rest energy in $E$. The Penning
trap, being a dynamics experiment (a charge in orbit), therefore reads energy \emph{as} inertia; it forces no mass
upon us, and RealQM, whose stationary states run no dynamics of their own (the timeless atom of the companion notes
on real time dependence~\cite{Time,Dynamics}), need never leave energy at all. Mass is primitive at neither stage: energy suffices for
what \emph{is}, and dresses as inertia only for what \emph{moves}.

This is not idle: it is why accelerators are built as they are. A cyclotron of fixed field and fixed
radio-frequency works only while the charge is slow, for its orbital frequency $\omega_c=qBc^2/E$ is constant only
while $E$ is near the rest energy. Speed the charge up and $E$ grows --- the invested energy \emph{becoming}
inertia --- so $\omega_c$ falls, the charge arrives late for each accelerating kick and slips out of phase, and the
machine stalls. The remedy (synchrocyclotron, then synchrotron) is to \emph{track the energy}: ramp frequency and
field to follow $\omega_c=qBc^2/E$ as $E$ climbs. An accelerator must chase energy, not mass; that it works is a
standing confirmation that the inertia it fights is $E/c^2$, growing without bound as $v\to c$ --- which is why $c$
is never reached, the energy going ever more into inertia and ever less into speed.

\begin{thebibliography}{9}
\bibitem{RealQM} C.~Johnson, \emph{Real Quantum Mechanics} (book).
\bibitem{Abraham} M.~Abraham, ``Prinzipien der Dynamik des Elektrons,''
  \emph{Ann.\ Phys.} \textbf{10}, 105--179 (1903).
\bibitem{Thomson} J.~J.~Thomson, ``On the electric and magnetic effects produced by the motion
  of electrified bodies,'' \emph{Philos.\ Mag.} \textbf{11}, 229--249 (1881).
\bibitem{Madelung} E.~Madelung, ``Quantentheorie in hydrodynamischer Form,''
  \emph{Z.\ Phys.} \textbf{40}, 322--326 (1927).
\bibitem{Bohm} D.~Bohm, ``A suggested interpretation of the quantum theory in terms of
  `hidden' variables. I,'' \emph{Phys.\ Rev.} \textbf{85}, 166--179 (1952).
\bibitem{Rainville} S.~Rainville \emph{et al.}, ``A direct test of $E=mc^2$,''
  \emph{Nature} \textbf{438}, 1096--1097 (2005).
\bibitem{Dehmelt} H.~Dehmelt, ``Experiments with an isolated subatomic particle at rest,''
  \emph{Rev.\ Mod.\ Phys.} \textbf{62}, 525--530 (1990).
\bibitem{Time} C.~Johnson, \emph{Real Time Dependence in RealQM: the Timeless Atom},
  RealMolecule gallery.
\bibitem{Dynamics} C.~Johnson, \emph{The Atom Does Not Run Its Own Dynamics},
  RealMolecule gallery.
\end{thebibliography}

\end{document}
